\documentclass[11pt]{article}

% ── Encoding & fonts ──────────────────────────────────────────────────────────
\usepackage[T1]{fontenc}
\usepackage[utf8]{inputenc}
\usepackage{lmodern}
\usepackage{microtype}

% ── Page geometry ─────────────────────────────────────────────────────────────
\usepackage[margin=1.25in]{geometry}

% ── Mathematics ───────────────────────────────────────────────────────────────
\usepackage{amsmath}
\usepackage{amssymb}
\usepackage{amsthm}

% ── Tables ────────────────────────────────────────────────────────────────────
\usepackage{booktabs}
\usepackage{array}
\usepackage{tabularx}

% ── Lists ─────────────────────────────────────────────────────────────────────
\usepackage{enumitem}
\setlist{itemsep=2pt, topsep=4pt}

% ── Hyperlinks ────────────────────────────────────────────────────────────────
\usepackage[colorlinks=true, linkcolor=black, citecolor=black,
            urlcolor=black, pdfborder={0 0 0}]{hyperref}
\usepackage{url}

% ── Bibliography ──────────────────────────────────────────────────────────────
\usepackage[numbers, sort&compress]{natbib}
\bibliographystyle{unsrtnat}

% ── Paragraph spacing ─────────────────────────────────────────────────────────
\usepackage{parskip}
\setlength{\parskip}{6pt}

% ── Section formatting ────────────────────────────────────────────────────────
\usepackage{titlesec}
\titleformat{\section}{\large\bfseries}{\thesection.}{0.6em}{}
\titleformat{\subsection}{\normalsize\bfseries}{\thesubsection}{0.6em}{}
\titlespacing{\section}{0pt}{14pt}{6pt}
\titlespacing{\subsection}{0pt}{10pt}{4pt}

% ── Abstract ──────────────────────────────────────────────────────────────────
\usepackage{abstract}
\renewcommand{\abstractnamefont}{\normalfont\bfseries}
\renewcommand{\abstracttextfont}{\normalfont\small}

% ── Footnotes ─────────────────────────────────────────────────────────────────
\usepackage[hang, flushmargin]{footmisc}

% ── Theorem-like environments ─────────────────────────────────────────────────
\theoremstyle{definition}
\newtheorem{definition}{Definition}
\theoremstyle{plain}
\newtheorem*{informalclaim}{Claim (informal)}
\newtheorem*{property}{Property}

% ─────────────────────────────────────────────────────────────────────────────
\title{\textbf{The Capture Problem:\\
Why Epistemic Decentralisation Is Not Sufficient,\\
and What Capture Resistance Requires}}

\author{%
  Dr.\ Richard Nicholson\\[4pt]
  \small Tathata Systems Ltd, United Kingdom\\
  \small \href{mailto:tathatasystems@proton.me}{tathatasystems@proton.me}
}

\date{June 2026\\[4pt]
\small \textbf{Working draft.} Companion to~\cite{nicholson2026,nicholson2026hlks}.}
% ─────────────────────────────────────────────────────────────────────────────

\begin{document}

\maketitle

\begin{abstract}
The companion papers establish that coordinator-based architectures fail
structurally~\cite{nicholson2026} and that coordinator-free, locally-resolved
substrate design is the unique \emph{epistemically} correct response across four
independent domains~\cite{nicholson2026hlks}. This paper argues that epistemic
correctness is \emph{necessary but not sufficient}. A substrate that resolves
knowledge locally---that satisfies the four correctness properties of the
companion paper---can nonetheless be \emph{captured}: its emergent consensus
mechanisms can be seized by concentrated interests, or an internal coordination
class can entrench and extract coordination rent, without any violation of
epistemic correctness. Capture is a failure of \emph{power} distribution, not of
\emph{knowledge} distribution, and the two are orthogonal. We introduce four
further substrate properties---Capture Resistance, Mandate TTL, Epistemic
Symmetry, and Exit---that an epistemically-correct substrate must additionally
satisfy to resist capture, and argue, following the cross-domain method of the
companion paper, that these properties appear independently in the design of
capture-resistant monetary systems (Bitcoin), capture-resistant governance
(Ostrom's commons, recallable-mandate confederalism), and capture-resistant
computing substrates. We relocate the MCB/P/S evaluative framework (Moral Circle
Breadth, Polycentricity, Surplus-Claim Regime) to its proper home here, as a
scorecard for capture resistance. We then confront the hardest question the
research programme faces: \emph{does coordinator-freeness prevent capture, or
merely relocate the captor to a place that is harder to see?} We argue---against a
common reading---that Bitcoin's surface concentration of mining, development, and
wealth does not constitute capture: its value resides in holders who are motivated
to resist seizure and able to exit, so that concentration is the \emph{source} of an
emergent immune response rather than evidence of capture, and a would-be captor
inherits a worthless network. Bitcoin shows capture resistance can be emergent and
robust; the sharper open question is whether a \emph{computing} substrate---which
lacks an internal asset whose value participants are individually motivated to
protect---can manufacture the same alignment, or must rely on weaker mechanisms. We
close with an honest account of what is open, including whether the present authors'
own substrate can supply that alignment or only makes its coordination class visible
and contestable.
\end{abstract}

% ─────────────────────────────────────────────────────────────────────────────
\section{Introduction}

The companion paper~\cite{nicholson2026hlks} establishes a strong result: across
economics, governance, organisations, and distributed computing, coordinator-free
locally-resolved substrate design is the unique structurally correct response to
the \emph{epistemic} problem of heterogeneous local knowledge. Four substrate
correctness properties---local signal propagation, emergent group formation,
minimal federation invariants, and subsidiarity-respecting consensus---follow from
the Coordinator Trap Theorem and appear, independently derived, in each domain.

That result is about \emph{knowledge}. This paper is about \emph{power}. Its thesis
is that the two are orthogonal, and that the first does not buy the second:

\begin{informalclaim}[The Insufficiency of Epistemic Correctness]
A substrate may satisfy all four epistemic-correctness properties of the companion
paper and remain vulnerable to capture. Capture is the seizure of a system's
coordination surplus by a subset of its participants or by external interests;
it is a failure of power distribution, and no property concerning the distribution
of \emph{knowledge} precludes it. Resistance to capture is therefore a distinct
structural axis, requiring its own properties.
\end{informalclaim}

The intuition is direct. The companion paper's Property~4 admits an emergent
consensus mechanism for the narrow class of operations requiring binding
cross-agent commitment. That mechanism is exactly the surface a captor seeks: it
is where the system's collective decisions are made, and a participant who acquires
disproportionate influence over it can bend system-level outcomes to extract rent
without ever telling a lie about local state. The substrate remains epistemically
correct---every agent still resolves its own knowledge locally---while becoming, at
the level that matters, an instrument of the captor. Epistemic correctness is the
floor, not the ceiling.

\paragraph{Contributions.} This paper:
\begin{enumerate}
\item States the orthogonality of epistemic correctness and capture resistance,
  and the two failure modes---\emph{external} and \emph{internal} capture---that
  an epistemically-correct substrate remains exposed to (\S\ref{sec:problem}).
\item Defines four further substrate properties---Capture Resistance, Mandate TTL,
  Epistemic Symmetry, and Exit---and argues that an epistemically-correct substrate
  must satisfy them to resist capture (\S\ref{sec:properties}).
\item Relocates the MCB/P/S framework to its proper role as a capture-resistance
  scorecard, and shows the three axes measure exactly the distinctions the
  capture-resistance properties imply (\S\ref{sec:mcbps}).
\item Instantiates the properties across domains---capture-resistant money
  (Bitcoin), capture-resistant governance (Ostrom; recallable-mandate
  confederalism), and capture-resistant computing---by the same convergence method
  as the companion paper (\S\ref{sec:domains}).
\item Confronts the central difficulty: whether coordinator-freeness prevents
  capture or relocates it, using Bitcoin's empirical concentration as the test
  case, and applies the same scrutiny to the present authors' own substrate
  (\S\ref{sec:hard}).
\item Gives an honest accounting of what remains open (\S\ref{sec:open}).
\end{enumerate}

\paragraph{Scope of the claim.} As with the companion paper, the cross-domain
argument is one of \emph{convergent independent derivation}, not mechanism
transfer: a capture-resistant computing substrate implements neither Bitcoin's
proof-of-work nor a confederal council, but exhibits the same structural
\emph{properties} that those independently arrived at. The pattern recurs; the
mechanisms do not transfer. Crucially, the structural analysis of an example is
\emph{orthogonal to endorsing it}: we analyse the capture-resistance structure of
Bitcoin and of recallable-mandate governance without endorsing the monetary
politics of the one or the revolutionary politics of the other. This separation is
not evasion---it is the precise discipline the cross-domain method requires, and
the place where the structural content of a charged example must be taken
seriously rather than avoided. Capture is, unavoidably, a normative-adjacent
subject (it concerns rent, power, and who a system serves); we keep the
\emph{structural} claim---which substrate designs make capture costly---separable
throughout from the \emph{normative} question of which distributions of power are
just.

% ─────────────────────────────────────────────────────────────────────────────
\section{The Capture Problem}
\label{sec:problem}

\subsection{What Capture Is}

\begin{definition}[Capture]
A system is \emph{captured} when a subset of its participants, or an external
interest, acquires disproportionate control over the system's binding coordination
decisions and uses that control to redirect the system's coordination surplus
toward itself. Capture is distinguished from ordinary influence by two features:
the control is \emph{disproportionate} to the captor's legitimate stake, and it
is used to extract \emph{rent}---value not matched by contribution.
\end{definition}

Capture presupposes a surface to capture. In a fully local system with no binding
collective decisions there is nothing to seize. But the companion paper shows that
some operations---those requiring linearisable cross-agent commitment---cannot be
resolved locally and require an emergent consensus mechanism (Property~4). That
mechanism is the capture surface. Capture resistance is therefore not the absence
of collective decision-making but the structuring of it so that disproportionate
control is costly to acquire and costly to retain.

\subsection{Two Failure Modes}

\paragraph{External capture.} A concentrated outside interest acquires control of
the emergent consensus mechanism and bends it to extract rent. The engine is the
asymmetry Olson~\cite{olson1965} identified: the benefits of capturing a coordination
mechanism are concentrated on the captor, while the costs are diffused across all
participants, so the captor is motivated to invest in capture far beyond any
individual defender's motivation to resist. The cost of capture and the rent it
yields are not symmetric; substrate design that does not engineer that asymmetry
in the defender's favour is, over a long enough horizon, captured by default.

\paragraph{Internal capture.} A subset of the system's \emph{own} participants---
those who operate, maintain, or hold the largest stake in the consensus
mechanism---entrenches and extracts. This is the subtler and, we will argue, the
harder mode: a system can be coordinator-\emph{free} in its protocol and still grow
an internal \emph{coordination class}---the maintainers who set the defaults, the
large stakeholders whose preferences dominate emergent consensus, the operators of
the points where federation seams meet. Internal capture does not require anyone to
violate a protocol rule. It requires only that influence over the collective
mechanism concentrate faster than the mechanism redistributes it.

These two modes correspond to the two failure modes visible in the empirical record
and named in the companion paper's open questions: acute seizure from without, and
chronic erosion from within.

% ─────────────────────────────────────────────────────────────────────────────
\section{Four Properties of a Capture-Resistant Substrate}
\label{sec:properties}

We state four properties additional to the companion paper's Properties 1--4. As
there, they are stated as structural conditions and argued to appear independently
across domains; their formal status is the same as the companion paper's---argued
with care, not proved in a single cross-domain formalism (\S\ref{sec:open}).

\begin{property}[5: Capture Resistance]
The cost $C_{\text{capture}}$ of acquiring disproportionate control of the
substrate's binding-commitment mechanism is high relative to the rent
$R_{\text{capture}}$ that control yields, and rises with the value the substrate
coordinates. A substrate is capture-resistant to the degree that
$C_{\text{capture}} / R_{\text{capture}}$ is large and increasing in stake.
\end{property}

The canonical engineered instance is Bitcoin's proof-of-work: control of the
consensus requires acquiring a majority of a deliberately expensive resource, and
the cost rises with the value secured. The computing-substrate analogue is open
and discussed in \S\ref{sec:hard}; it is \emph{not} simply ``make consensus
expensive,'' because gratuitous cost violates the companion paper's
efficiency-by-locality commitments. The design tension---capture resistance
demands that control be costly, epistemic correctness demands that coordination be
cheap---is real and is a central subject of this paper.

\begin{property}[6: Mandate TTL]
No grant of influence over the collective mechanism is permanent. Authority is
held under a time-to-live and must be actively re-granted; absent renewal, it
evaporates. Capture that depends on holding a position is bounded by the rate at
which positions must be re-earned.
\end{property}

Mandate TTL is the governance analogue of the companion substrate's defining
mechanism: capability advertisements, group memberships, and load signals are all
soft-state under a TTL, and a node that stops re-advertising simply fades. The
companion paper applies this to \emph{liveness}; here it applies to
\emph{authority}. Recallable delegation, fixed terms, and graduated re-grant are
the governance instances; ``management as evaporating intent''---an instruction
that must be continuously re-asserted to persist, and that any node may locally
override---is the computing instance. A permanent mandate is a standing capture
surface; a mandate that must be re-earned each interval is not.

\begin{property}[7: Epistemic Symmetry]
No participant holds privileged access to the substrate's state that other
participants structurally lack. Where one party can observe the global state and
others cannot, that asymmetry is itself a capture vector: the observer can act on
information unavailable to those it coordinates.
\end{property}

Epistemic symmetry connects capture resistance back to the data-locality argument
of the companion paper, and to a result of independent practical importance: a
substrate in which \emph{data stays where it is produced}---at the edge, under the
control of the participant that generated it---denies any central party the
asymmetric view that capture exploits. Data sovereignty, which the engineering
literature treats as a compliance property, is here a \emph{capture-resistance}
property: the same structural feature that keeps a regulated participant's data
under its own control also denies a would-be captor the privileged vantage point.

\begin{property}[8: Exit]
A participant, or a coherent sub-domain of participants, can unilaterally withdraw
from the substrate---taking its state and its capabilities with it---without the
permission of any other party. Exit is the capture-resistance mechanism that does
not depend on the captor playing fair.
\end{property}

Exit is Hirschman's~\cite{hirschman1970} insight applied to substrate design. Voice
(participating in the collective mechanism) and loyalty (remaining despite
dissatisfaction) both leave the captor in place; only exit removes the captor's
hold. A federation of sovereign sub-domains, each able to leave at will, bounds
capture at the federation level structurally: no federation-wide captor can hold a
sub-domain that can simply depart. Exit is the property most directly served by the
companion substrate's sovereign federation model, and the one that most sharply
distinguishes a federation of equals from a hierarchy with extra steps.

% ─────────────────────────────────────────────────────────────────────────────
\section{The MCB/P/S Framework as a Capture Scorecard}
\label{sec:mcbps}

The MCB/P/S framework---Moral Circle Breadth, Polycentricity, Surplus-Claim
Regime---was introduced by the present author~\cite{nicholson2025} for the
evaluation of monetary architectures, and borrowed by the companion
paper~\cite{nicholson2026hlks} to score epistemic distinctions. As that paper
notes, its proper home is the analysis of capture, which we develop here.

\begin{itemize}
\item \emph{Moral Circle Breadth} measures \emph{who the substrate serves}---whose
  interests are inside the boundary the system optimises for. A captured substrate
  has a narrow effective moral circle (the captor) regardless of its nominal
  membership.
\item \emph{Polycentricity} measures \emph{how distributed binding decision-making
  is}---the fraction of decisions resolved locally versus through a single
  collective mechanism. Low polycentricity is a large capture surface.
\item \emph{Surplus-Claim Regime} measures \emph{who claims the coordination
  surplus}---the rent the system generates. This is the capture axis proper: a
  captured substrate has a concentrated surplus-claim regime.
\end{itemize}

The three axes are not independent of Properties 5--8; they are the observable
\emph{symptoms} the properties are designed to prevent. A substrate scoring well on
Capture Resistance, Mandate TTL, Epistemic Symmetry, and Exit will, we argue, score
broad-MCB, high-P, distributed-surplus; and the failure of any property shows up as
degradation on a corresponding axis. Table~\ref{tab:propaxis} states the mapping.

\begin{table}[t]
\centering
\small
\renewcommand{\arraystretch}{1.2}
\begin{tabularx}{\textwidth}{@{}llX@{}}
\toprule
\textbf{Property} & \textbf{Primary axis} & \textbf{Symptom if the property fails} \\
\midrule
5 \,Capture Resistance & Surplus-Claim & Low $C_{\text{capture}}$ lets a captor seize the consensus mechanism and claim coordination rent; surplus concentrates. \\
6 \,Mandate TTL & Polycentricity & Authority accretes into durable positions; binding decisions centralise into a standing class. \\
7 \,Epistemic Symmetry & Moral Circle Breadth & A privileged observer acts on a view others lack; the effective circle narrows to that observer's interests. \\
8 \,Exit & Surplus-Claim / MCB & Captive participants cannot leave; extractable rent is unbounded and the served circle shrinks to those who cannot exit. \\
\bottomrule
\end{tabularx}
\caption{Each capture-resistance property protects a primary MCB/P/S axis; failure
of the property is observable as degradation on that axis. The MCB/P/S framework is
thus the \emph{scorecard} for the properties, not an independent construct.}
\label{tab:propaxis}
\end{table}

% ─────────────────────────────────────────────────────────────────────────────
\section{Domain Instantiations of Capture Resistance}
\label{sec:domains}

By the convergence method of the companion paper, we expect the capture-resistance
properties to appear independently wherever capture has been a live design concern.

\subsection{Capture-Resistant Money: Bitcoin}

Bitcoin~\cite{nakamoto2008} is the first \emph{engineered} capture-resistant
substrate, and it instantiates capture resistance more richly than any single
mechanism suggests: an engineered cost-of-capture at the protocol layer, and---above
it---an emergent economic immune response that a surface reading routinely mistakes
for its opposite.

At the protocol level its capture resistance is strong and deliberate. Control of
the consensus requires a majority of a resource (hash power) made costly by design,
with a cost that rises as the value secured rises---an explicit, scaling
$C_{\text{capture}}$ (Property~5). There is no issuer, foundation, or authority
empowered to alter the monetary schedule, so control of the protocol carries no
permanent mandate (Property~6). The ledger is public and verified identically by
every full node, so no participant holds a privileged view of system state
(Property~7). And a participant who rejects a ruleset may fork, carrying its state
into a substrate governed by different rules (Property~8). On the four properties,
the protocol scores well, and did so by explicit design rather than by accident.

At the emergent level its capture resistance is---contrary to a common reading---
also strong, though by a different mechanism than the protocol provides, and the
distinction is this paper's central lesson about where capture resistance actually
lives. The substrate does exhibit surface concentration in the layers the protocol
does not govern: hash power concentrates into a few mining pools, development into a
small set of maintainers, and holdings into a small fraction of addresses. The naive
inference---that this concentration constitutes capture---is wrong, and seeing why
is the point.

\emph{Concentration of a layer is not capture of the system, because capture
requires control of the layer that holds value, and in Bitcoin that layer is the one
most resistant to seizure.} A mining coalition, even a majority one, cannot rewrite
the monetary rules, steal balances, or impose a ruleset the economic majority
rejects. The value of the network resides in its holders and in the social consensus
over which chain is legitimate; any coalition that attacks the rules destroys the
value it would need to profit from the attack. The attack is self-defeating. The
parties with the most to lose from capture---the holders of the majority of the
network's value---are therefore not its captors but its immune system: they have the
strongest motivation to neutralise a captor and the decisive means to do so, because
they can exit (sell, or fork to a chain that excludes the attacker) and value follows
the legitimate chain rather than the attacker's. Property~8 (Exit) is doing the work,
at the economic rather than the protocol layer. (Concentration of hash power can
still enable bounded, detectable harms short of capture---transaction censorship, a
short reorganisation---but even these are disciplined by the same response: a pool
that misbehaves bleeds the hash power it rents, because the miners who supply it
answer to the same value signal.)

This is the elegant part, and the part a surface reading misses. Bitcoin's capture
resistance is not a single engineered barrier but a multi-layered equilibrium: an
engineered cost-of-capture at the protocol layer, and above it an emergent immune
response in which the \emph{concentration of value}, far from being the vulnerability,
is the \emph{source} of resistance, because the concentrated holders are precisely
those most motivated to defeat a captor. The 2017 dispute over block size is the
demonstration: a coalition of the most powerful concentrated parties---large
businesses and a mining majority---attempted a rule change, and the economic users
and nodes rejected it; the change was abandoned. Every surface metric of
concentration was present, and capture still failed. Concentration and capture are
different things, and conflating them is the error this paper most wants to avoid.

\subsection{Capture-Resistant Governance: Ostrom and Recallable Mandate}

Read structurally rather than normatively, Ostrom's~\cite{ostrom1990}
empirically-derived design principles for long-enduring commons are a catalogue of
capture-resistance mechanisms. Graduated sanctions and accountable monitoring raise
the cost of defecting from or seizing the collective-choice arrangement
(Property~5). Nested enterprises decompose authority so that no single tier holds
decisive control, keeping decisions polycentric (Property~6 at each tier). The
principle that appropriators may modify their own operating rules keeps mandate
local and revocable rather than imposed and permanent (Property~6). And clearly
defined boundaries, paired with the practical ability of appropriators to abandon a
failing arrangement, supply a form of Exit (Property~8). Ostrom did not derive these
from a theory of capture; she observed which commons endured for centuries and which
collapsed, and the enduring ones independently exhibit the capture-resistance
properties. That is the convergence this paper claims.

The sharpest instantiation of Mandate~TTL (Property~6) is found in
\emph{recallable-mandate confederalism}, of which the contemporary governance of
Rojava (the Autonomous Administration of North and East
Syria)~\cite{ocalan2011,knapp2016} is the most fully realised modern instance.
Authority is delegated upward through councils under \emph{continuous recall}:
a delegate's mandate may be withdrawn at any time by those who granted it, and
binding decisions originate in local assemblies rather than descending from a
centre. Structurally, this is Mandate~TTL taken to its limit---authority that is
not merely term-limited but revocable at every instant---and it is paired with
strong Polycentricity (decisions default to the most local competent assembly) and
a deliberately broad Moral Circle (the boundary principle is inclusion by
participation, not by prior status). We analyse this structure for its
capture-resistance content alone. As stated in the introduction's scope of the
claim, the structural claim is orthogonal to the political programme in which the
structure is embedded: that continuously revocable mandate raises the cost of durable capture is
a claim about \emph{authority dynamics}, not an endorsement of a revolutionary
movement, and it survives substitution by any far less charged arrangement that
makes mandate continuously revocable---frequent recall elections, mandatory
rotation, sunset clauses on delegated powers. We name Rojava not for its politics
but because it is where the structure is implemented most completely, and a
cross-domain argument that quietly omitted its clearest governance instance to
avoid discomfort would be failing its own method.

\subsection{Capture-Resistant Computing: the Companion Substrate}

The companion substrate instantiates the properties partially and unevenly, and
honesty requires saying which.

It provides \emph{Exit} structurally (Property~8): a federation of self-certifying
sovereign sub-domains, each able to depart unilaterally---carrying its own state and
identity---without the permission of any other party, because no issuing authority
exists whose consent departure would require. It provides \emph{Epistemic Symmetry}
(Property~7) through data locality---state is held by the participant that produces
it and replicated only as the soft-state its consumers need---and through
self-certified identity, so that no party holds an issuing key conferring a
privileged trust position. It provides a form of \emph{Mandate~TTL} (Property~6)
through evaporating soft-state and management-as-intent: an operator instruction is
expressed as an intent that must be continuously re-asserted to persist and that any
node may locally override, so no operator mandate becomes permanent.

Its \emph{Capture Resistance} (Property~5) is the least developed, and the
substrate's posture here is \emph{detection rather than prevention}. It has no
engineered cost-of-capture analogous to proof-of-work; instead it makes capture
\emph{visible}. A hash-chained, signed transparency log records changes to the
trust set, so that a captor cannot quietly rewrite who is trusted; namespace-
ownership tripwires flag writes that cross authority boundaries; and an ongoing
diagnosability effort aims to make emergent, fleet-level behaviour---including the
formation of an internal coordination class---legible to a non-privileged observer.
Detection is structurally weaker than prevention: it does not stop capture, it
denies the captor the ability to act \emph{quietly}, raising the cost of capture by
the cost of being seen. Whether, for a computing substrate, detection-plus-exit is
an adequate substitute for an engineered $C_{\text{capture}}$---or whether such
substrates need a prevention mechanism of their own---is the open question carried
into \S\ref{sec:hard} and \S\ref{sec:open}.

% ─────────────────────────────────────────────────────────────────────────────
\section{The Hard Question: Prevention or Relocation?}
\label{sec:hard}

We state the difficulty the programme cannot avoid:

\begin{quote}
\emph{Does coordinator-freeness prevent capture, or merely relocate the captor to a
place that is harder to see?}
\end{quote}

Bitcoin is the test case that sharpens the question---but not in the direction a
surface reading suggests. As \S\ref{sec:domains} argued, Bitcoin's surface
concentration does not constitute capture: its value layer resists seizure because
the holders are motivated to resist and able to exit, so a would-be captor inherits
a worthless network. Bitcoin is therefore evidence that capture resistance can be
\emph{emergent and robust}, not merely engineered---encouraging for the programme
rather than damning. The captor did not relocate to an unbounded emergent class; it
was defeated by an emergent immune response the surface metrics do not show.

That result sharpens a harder and more specific question for a \emph{computing}
substrate, because Bitcoin's immune response rests on a feature such a substrate may
lack: an asset whose value every participant is individually motivated to protect.
In Bitcoin, everyone with holdings has skin in the game---an attack on the rules is
an attack on their wealth, and self-interest aligns the broad majority against any
captor automatically. A substrate that coordinates AI agents has no equivalent stake
by default: its participants run workloads, but they do not necessarily hold a
fungible, exit-portable claim on the substrate's continued integrity whose value
would collapse if the substrate were captured. The disciplining value signal that
makes Bitcoin's exit decisive may simply be absent.

The honest question is therefore not \emph{whether} coordinator-freeness can prevent
capture in the abstract---Bitcoin shows the properties can deliver it---but whether a
computing substrate can \emph{manufacture the alignment} that makes capture
resistance emergent, or must instead rely on the weaker mechanisms available to it:
detection without a value signal, and exit without an asset whose flight punishes the
captor. The load in resisting internal capture still falls on Properties~6--8---a
continuously re-earned mandate (6) denies a class a durable position, epistemic
symmetry (7) denies it an informational edge, and exit (8) caps the rent it can
extract---but Bitcoin shows these are strongest when an aligned value signal stands
behind them. Whether a substrate can supply that signal without an internal asset,
or resist a coordination class without it, is, we believe, the central open problem
of coordinator-free design, and the one on which the normative significance of the
entire programme rests.

We apply the same scrutiny to our own substrate. A coordinator-free computing
substrate is not immune: its maintainers set defaults, its largest deployments
shape its emergent norms, and the operators of its federation seams occupy
structurally privileged positions. We claim only that it makes these positions
\emph{visible and contestable}---through transparency, mandate-TTL, and exit---not
that it abolishes them. A research programme that claimed to have abolished internal
capture would be making exactly the totalising error its own thesis warns against.

% ─────────────────────────────────────────────────────────────────────────────
\section{Open Questions}
\label{sec:open}

\begin{enumerate}
\item \emph{The formal status of capture resistance.} As with the companion paper's
  theorem, Properties 5--8 are argued informally to preserve cross-domain
  generality. Whether $C_{\text{capture}}$ and the surplus-claim regime admit a
  cross-domain quantification on a common scale is open.

\item \emph{The cost/correctness tension.} Capture resistance wants binding decisions
  to be costly to control; epistemic correctness wants coordination to be cheap.
  Whether a single substrate can satisfy both, and where the frontier lies, is open.

\item \emph{The sufficiency of Properties 6--8 against internal capture.} The hard
  question of \S\ref{sec:hard}: can mandate-TTL, epistemic symmetry, and exit bound
  an emergent coordination class, or does internal capture require a mechanism not
  yet identified?

\item \emph{An engineered $C_{\text{capture}}$ for computing.} Whether a
  coordinator-free computing substrate needs an explicit cost-of-capture analogous
  to proof-of-work, or whether detection-plus-exit suffices, is open and bears
  directly on the companion substrate's design.

\item \emph{The normative bridge.} This paper keeps the structural claim separable
  from the normative one. A complete treatment must eventually connect them: to say
  which \emph{distributions} of power a capture-resistant substrate tends to
  produce, and whether structural capture resistance reliably yields outcomes a
  normative theory would endorse---or merely makes capture harder without making it
  rare.
\end{enumerate}

% ─────────────────────────────────────────────────────────────────────────────
\bibliography{references}

\end{document}
